\begin{frame}{Código Gray: Por que utilizar?}
    \small

    Considere a transição do número 3 para o número 4.

    \begin{columns}[T]

        \begin{column}{0.48\textwidth}

            \begin{block}{Binário}
                \[
                    3 = 011
                \]

                \[
                    4 = 100
                \]

                Portanto:

                \[
                    011 \rightarrow 100
                \]

                \textbf{Três bits mudam.}
            \end{block}

        \end{column}

        \begin{column}{0.48\textwidth}

            \begin{block}{Gray}
                \[
                    3 = 010
                \]

                \[
                    4 = 110
                \]

                Portanto:

                \[
                    010 \rightarrow 110
                \]

                \textbf{Apenas um bit muda.}
            \end{block}

        \end{column}

    \end{columns}

    \vspace{0.3cm}

    \begin{alertblock}{Vantagem}
        Isso reduz erros durante transições entre estados.
    \end{alertblock}

\end{frame}